\documentclass[../zusammenfassung_AN.tex]{subfiles}
\begin{document}
	
	\subsection{Grundfunktionen im rechtwinkligen Dreieck}
	
	\begin{tcolorbox}[dglbox, title={Definitionen (G = Gegenkathete, A = Ankathete, H = Hypotenuse)}]
		\[
		\sin(\alpha)=\frac{G}{H} \qquad
		\cos(\alpha)=\frac{A}{H} \qquad
		\tan(\alpha)=\frac{G}{A}=\frac{\sin(\alpha)}{\cos(\alpha)} \qquad
		\cot(\alpha)=\frac{A}{G}=\frac{\cos(\alpha)}{\sin(\alpha)}=\frac{1}{\tan(\alpha)}
		\]
	\end{tcolorbox}
	
	% -------------------------------------------------------
	\subsection{Sinus- und Kosinussatz}
	% -------------------------------------------------------
	
	\begin{tcolorbox}[formelbox]
		\textbf{Sinussatz:}
		\[
		\frac{a}{\sin(\alpha)} = \frac{b}{\sin(\beta)} = \frac{c}{\sin(\gamma)}
		\]
		\textbf{Kosinussatz:}
		\[
		a^2 = b^2+c^2-2bc\cos(\alpha) \qquad
		b^2 = a^2+c^2-2ac\cos(\beta) \qquad
		c^2 = a^2+b^2-2ab\cos(\gamma)
		\]
	\end{tcolorbox}
	
	\begin{tcolorbox}[hinweis]
		\textbf{SSW-Fall (Sinussatz):} Es können zwei Lösungen auftreten, da
		$\sin(\beta)=\sin(180^\circ-\beta)$.
	\end{tcolorbox}
	
	% -------------------------------------------------------
	\subsection{Trigonometrische Identitäten}
	\label{subsubsec:trigonometrische_identitäten}
	% -------------------------------------------------------
	
	\subsubsection{Beziehungen zwischen Sinus und Kosinus}
	
	\begin{tcolorbox}[formelbox]
		\[
		\sin(x)=\cos\!\left(x-\tfrac{\pi}{2}\right) \qquad
		\cos(x)=\sin\!\left(x+\tfrac{\pi}{2}\right) \qquad
		{-\sin(x)}=\cos\!\left(x+\tfrac{\pi}{2}\right) \qquad
		{-\cos(x)}=\sin\!\left(x-\tfrac{\pi}{2}\right)
		\]
	\end{tcolorbox}
	
	\subsubsection{Pythagoräische Identität und Symmetrie}
	
	\begin{tcolorbox}[formelbox]
		\[
		\sin^2(\alpha)+\cos^2(\alpha)=1
		\qquad
		\cos(x)=\cos(-x)
		\qquad
		\sin(x)=-\sin(-x)
		\]
	\end{tcolorbox}
	
	\subsubsection{Additionstheoreme}
	
	\begin{tcolorbox}[dglbox, title={Additionstheoreme}]
		\textbf{Sinus:}
		\[
		\sin(\alpha\pm\beta) = \sin(\alpha)\cos(\beta) \pm \sin(\beta)\cos(\alpha)
		\]
		\textbf{Kosinus:}
		\[
		\cos(\alpha+\beta) = \cos(\alpha)\cos(\beta)-\sin(\alpha)\sin(\beta) \qquad
		\cos(\alpha-\beta) = \cos(\alpha)\cos(\beta)+\sin(\alpha)\sin(\beta)
		\]
		\textbf{Tangens:}
		\[
		\tan(\alpha+\beta) = \frac{\tan(\alpha)+\tan(\beta)}{1-\tan(\alpha)\tan(\beta)} \qquad
		\tan(\alpha-\beta) = \frac{\tan(\alpha)-\tan(\beta)}{1+\tan(\alpha)\tan(\beta)}
		\]
		\textbf{Kotangens:}
		\[
		\cot(\alpha+\beta)=\frac{\cot(\alpha)\cot(\beta)-1}{\cot(\beta)+\cot(\alpha)} \qquad
		\cot(\alpha-\beta)=\frac{\cot(\alpha)\cot(\beta)+1}{\cot(\beta)-\cot(\alpha)}
		\]
	\end{tcolorbox}
	
	\subsubsection{Doppelwinkeltheoreme}
	
	\begin{tcolorbox}[hinweis]
		Spezialfall der Additionstheoreme für $\alpha=\beta$:
	\end{tcolorbox}
	
	\begin{tcolorbox}[formelbox]
		\[
		\sin(2\alpha)=2\sin(\alpha)\cos(\alpha)
		\]
		\[
		\cos(2\alpha)=\cos^2(\alpha)-\sin^2(\alpha) = 2\cos^2(\alpha)-1 = 1-2\sin^2(\alpha)
		\]
		\[
		\tan(2\alpha)=\frac{2\tan(\alpha)}{1-\tan^2(\alpha)} \qquad
		\cot(2\alpha)=\frac{\cot^2(\alpha)-1}{2\cot(\alpha)}
		\]
	\end{tcolorbox}
	
	\subsubsection{Potenzen — nützlich für Integration}
	
	\begin{tcolorbox}[dglbox, title={Quadratische Ausdrücke (aus Doppelwinkeltheorem)}]
		\[
		\sin^2(\alpha) = \tfrac{1}{2}\bigl(1-\cos(2\alpha)\bigr)
		\qquad
		\cos^2(\alpha) = \tfrac{1}{2}\bigl(1+\cos(2\alpha)\bigr)
		\]
	\end{tcolorbox}
	
	\subsubsection{Produktformeln}
	
	\begin{tcolorbox}[formelbox]
		\[
		\sin x\cos y = \tfrac{1}{2}\bigl(\sin(x+y)+\sin(x-y)\bigr)
		\]
		\[
		\cos x\cos y = \tfrac{1}{2}\bigl(\cos(x+y)+\cos(x-y)\bigr)
		\]
		\[
		\sin x\sin y = \tfrac{1}{2}\bigl(\cos(x-y)-\cos(x+y)\bigr)
		\]
	\end{tcolorbox}
	
\end{document}